% -------------------------------------
%   Future Work
% -------------------------------------
\section{Future Work}
\label{sec:disc_futurework}

The results of this work indicate that, at MNIST scale, TM inference on the NEORV32 core is limited by the movement of model data rather than by clause evaluation. Based on these result, there are multiple different directions for future work to explore which might yield better performance benefits than the ones found here. Various different ideas for how these directions could be explored. The first direction is to target the memory bottleneck directly, by reducing the cost of transferring the model out of DDR2 (or possibly another type of external) or by widening the memory interface. The remaining directions look at workloads and platforms in which clause evaluation occupies a larger share of inference time, where processor-level acceleration like the CFU may remain worthwhile. Lastly on-chip training is discussed.

% -------------------------------------
%   DDR2 burst transfers
% -------------------------------------
\subsection{DDR2 Burst Transfers}
\label{sec:disc_burst}

The per-class preload reads the model from DDR2 one 32-bit word at a time, and Section~\ref{sec:disc_bottleneck} attributed its low effective transfer rate to the overhead of each individual AXI transaction rather than to the DDR2 device itself. With a burst transfer, the fixed overhead is paid once for a whole group of words rather than once per word, so the cost carried by each word is much lower. Replacing the individual word reads with burst transfers is therefore expected to reduce the preload time substantially \cite{Nolting2024Datasheet}. Since the preload accounts for approximately 91\% of per-class inference time, this is the single change with the greatest potential effect on total inference time. The preload already reads sequential addresses, so the access pattern is well suited to bursting; the work required lies in the bus interface and the transfer routine rather than in the access order. 

% -------------------------------------
%   Hybrid load strategy
% -------------------------------------
\subsection{Hybrid Load Strategy with Early Exit}
\label{sec:disc_hybrid}

The per-chunk and per-class strategies represent two extremes. The per-chunk strategy reads only the portion of the model that the early exit does not skip, but pays the latency of a scattered DDR2 read for every word. The per-class strategy reads sequentially, which opens up the possibility for burst transfers, but forfeits the early exit and therefore transfers the entire model unconditionally. A hybrid strategy could combine the advantages of both. Rather than reading one word at a time or preloading a full class, the model could be read in small bursts of several words, with clause evaluation proceeding over each burst before the next is requested. This would retain the efficiency of burst transfers while still allowing a clause to terminate early, so that the remaining bursts of a failed clause are never read. The size of the burst would set the balance between transfer efficiency and the granularity of the early exit, and identifying a suitable burst length would be the central question of this direction.

% -------------------------------------
%   Wider memory interface and external accelerator
% -------------------------------------
\subsection{Wider Memory Interface and a Processor-External Accelerator}
\label{sec:disc_external}

The constraint underlying the memory bottleneck is the width of the memory interface: the processor accesses DDR2 across the XBUS 32 bits at a time, which bounds how quickly the model can be moved regardless of the access strategy. Within a 32-bit core this width is fixed. By moving to a 64-bit architecture one would double the data moved per access, but the improvement would be bounded by roughly a factor of two which, while effective, does not change the memory-bound characteristic of the problem.

A more substantial gain would require moving model storage and clause evaluation outside the processor entirely. A processor-external accelerator, attached for example through the XBUS, is not constrained by the 32-bit register width of the CPU and could be given a wider and more numerous set of inputs. With the clauses of each class held in separate memory banks, the classes could then be evaluated in parallel rather than sequentially, addressing both the transfer width and the serial evaluation of the ten classes at once. The practical limit on such a design would likely be fan-in, as supplying a clause evaluation unit with as many literals as a full clause in parallel becomes difficult to route and time as the literal count grows. Determining a the best balance between parallelism and fan-in would therefore likely be the primary design challenge in this direction.

% -------------------------------------
%   Other TM variants
% -------------------------------------
\subsection{Acceleration of Other Tsetlin Machine Variants}
\label{sec:disc_variants}

The finding that clause evaluation is a small part of inference time is specific to the vanilla TM evaluated here, in which evaluation reduces to bitwise operations. Other TM variants, like the Convolutional TM (CTM) and the Coalesced TM (CoTM) \cite{GranmoOle-Christoffer2019TCTM,GlimsdalSondre2021CMTM}, have other characteristics which might mean they benefit more from acceleration through custom instructions.

The most important difference is that these variants reach the same accuracy with far fewer clauses. The CoTM shares clauses across classes by attaching a weight to each clause, so that a single clause can contribute to several classes at once \cite{GlimsdalSondre2021CMTM}. As a result it reaches an accuracy comparable to a vanilla TM while using only about 100 clauses per class, where the vanilla TM required several thousand, reducing the model size by more than two orders of magnitude \cite{GlimsdalSondre2021CMTM}. The CTM applies each clause across many small patches of an image rather than once over the whole image, which raises accuracy without a proportional increase in the number of clauses \cite{GranmoOle-Christoffer2019TCTM}. Since each patch is smaller than the full image, fewer literals are processed simultaenously, but this comes at the expense of longer processing time due to needing to process multiple different patches. Because the size of the model, and therefore the amount of data that must be moved out of DDR2, grows with the number of clauses, variants needing far fewer clauses produces a far smaller model and far fewer memory accesses. A sufficiently small model could even fit in on-chip memory, which would remove the memory bottleneck entirely and return the workload to be compute-bound, like in the preceding project thesis where the custom instruction was effective (Section~\ref{sec:disc_cfu}).

These variants also make the evaluation of each clause more demanding, which further increases the share of inference time spent on computation. The CTM must generate and handle the image patches that each clause is applied to \cite{GranmoOle-Christoffer2019TCTM}, and the weights of the CoTM introduce integer multiplication into the vote computation \cite{GlimsdalSondre2021CMTM}. Clause processing and evaluation therefore takes up a larger part of inference time than in the vanilla TM, which means there are more operations that can be offloaded to hardware through custom instructions or other approaches.

For the CTM and CoTM variants, acceleration through either the CFU or the CFS could be worthwhile, with the choice depending on how much computation is moved into hardware. The CFS, being a memory-mapped peripheral that can operate independently of the CPU (Section~\ref{sec:neorv32_cfs}), is better suited than the CFU to the larger and more stateful processing these variants might require \cite{Nolting2024Datasheet}.

% -------------------------------------
%   Existing ISA extensions
% -------------------------------------
\subsection{Existing ISA Extensions and Alternative Cores}
\label{sec:disc_vector}

The acceleration in this work was achieved through a custom instruction. An alternative is to use a standard ISA extension rather than a custom one, in particular the RISC-V vector (V) extension, which provides operations over vectors of data, including bitwise vector operations, and could process multiple literal words per instruction, maybe even whole clauses \cite{riscv_unpriv_isa}. The V extension is not currently available on the NEORV32 \cite{Nolting2024Datasheet}, so this direction would require the use of a different processor. Some open-source RISC-V cores support the vector extension, such as Viscuna, and could be a suitable target for such an investigation \cite{platzer2021vicuna}. 

Using a standardised extension rather than a custom one would also improve portability, as the resulting software would not depend on a processor-specific instruction. This approach could also be explored on platforms aside from FPGAs, as programmable logic would not really be the focus anymore, with boards like K230 providing a RISC-V processor supporting the V extension \cite{canaan_k230_datasheet}.

% -------------------------------------
%   On-chip training feasibility
% -------------------------------------
\subsection{On-Chip Training Feasibility}
\label{sec:disc_training}

As stated in Section \ref{sec:training_config}, on-chip training was excluded from this work, as the training procedure done on the host-pc relies on floating-point arithmetic that the NEORV32 executes only in software, in the absence of a floating-point unit (Section~\ref{sec:tm_learning}). This does not mean on-chip training is not possible, however it would be slow, and whatever performance benefits the CFU hopefully would provide would be overshadowed. In addition, the model would need to include TA states, which greatly increase the size of the model and increase the already large number of external memory reads. In other words, it is outside of the scope of the project.

A proof of concept at reduced scale, however, would likely be feasible. At a smaller configuration, for example a few hundred clauses per class trained on a reduced number of samples, the full TA state of all classes would fit within DDR2, and a training pass would complete in a time acceptable for demonstration rather than production use. Such an effort would establish whether the platform can support the complete TM life cycle of training and inference on-device, rather than relying on a model trained on a host. Whether custom instructions could meaningfully accelerate the training procedure, given its dependence on floating-point feedback, would be a further question for that investigation, but maybe some of the approaches suggested above could prove useful.